\chapter{LITERATURE REVIEW}
\label{chap:literature_review}

\section{Data acquisition and warehousing in e-commerce}

The first requirement of any market-intelligence system is reliable
access to marketplace data. In e-commerce contexts, relevant signals are
distributed across search-result pages, product-detail pages, image
assets, review pages, and advertising placements. Consequently, the
technical challenge is not limited to downloading web pages; it also
includes normalising heterogeneous records, preserving historical
snapshots, and maintaining a schema that supports downstream analytical
queries.

\subsection{From traditional scraping to actor-based extraction}

Traditional web scraping commonly relies on Document Object Model (DOM)
parsers, XPath selectors, CSS selectors, or browser automation tools
such as Selenium \citep{mitchell2018webscraping, kouzis2016scrapy}.
These techniques are useful for controlled and relatively stable web
pages, but they become fragile in modern e-commerce environments where
interfaces are dynamic, JavaScript-heavy, and frequently modified
through A/B testing. In addition, large retail platforms can make purely
custom scraping pipelines expensive to maintain because extraction logic
must be updated whenever page structure or access conditions change
\citep{mitchell2018webscraping, apify2024platform}.

Recent data-acquisition practice has therefore shifted toward
actor-based extraction services. Platforms such as Apify encapsulate
scraping logic in reusable cloud actors and return structured JSON
datasets for downstream processing \citep{apify2024platform}. This
approach does not remove the need for data validation, deduplication, or
schema design, but it reduces the engineering burden of maintaining
low-level browser automation. This thesis follows that approach: Amazon
category rankings, product snapshots, and reviews are collected through
Apify actors, while project-specific transformation and validation are
implemented in the Python ingestion scripts.

\subsection{Relational storage, orchestration, and reproducibility}

Once raw marketplace observations are collected, they must be stored in
a form that supports time-series comparison, metric calculation, and
agent queries. A relational warehouse is well suited to this task because
it enables historical snapshots, keyed upserts, and join-based retrieval
across products, categories, reviews, alerts, and derived metrics.
Supabase provides a managed PostgreSQL environment together with storage
and access-control features \citep{supabase2024docs}. In the implemented
system, Supabase tables store both raw records and analytical outputs,
while Supabase Storage is used for product-image persistence.

Scheduled orchestration is also central to reproducibility. GitHub
Actions supports scheduled workflows, manual dispatches, and repeatable
execution environments \citep{githubactions2024}. In this project, the
daily and weekly pipelines are implemented as scheduled GitHub Actions
jobs. This design is more appropriate for the present thesis than a
real-time streaming system, because the empirical objective is daily
market monitoring under cost and deployment constraints rather than
second-by-second surveillance.

\section{Commercial competitive intelligence and dashboard limitations}

Amazon sellers already have access to commercial intelligence platforms,
including tools such as Helium 10 and Jungle Scout
\citep{helium10listing, junglescout2024}. These systems provide useful
visibility into historical prices, product rankings, listing quality,
keyword competition, and product opportunities. They demonstrate that
there is real demand for marketplace intelligence and that sellers value
both quantitative signals, such as Best Sellers Rank (BSR), and
qualitative signals, such as review content \citep{amazon2024bsr,
chevalier2006, mudambi2010helpful}.

However, many commercial tools remain dashboard-centred. Dashboards are
valuable for exploration, but they require users to actively open the
interface, configure filters, inspect charts, and infer the operational
meaning of observed changes. This pull-based workflow creates cognitive
friction, especially when the user must monitor multiple categories,
products, and signals simultaneously. Earlier work in competitive
intelligence emphasises the managerial value of timely information
\citep{gilad1985}, while e-commerce and sponsored-search studies show
that reviews, ranking visibility, and paid placement can affect digital
market outcomes \citep{chevalier2006, ghose2009, mudambi2010helpful}.
Together, these perspectives suggest that the value of data depends not
only on collection, but also on how quickly and clearly it is transformed
into decisions.

For this reason, the present thesis positions dashboards as only one
delivery modality. The Streamlit and web dashboards support analytical
exploration, while the Telegram-based agent layer is designed to push
short, role-specific briefs when the system identifies relevant market
events.

\section{Machine learning for marketplace signal extraction}

Raw marketplace records become useful only after they are transformed
into interpretable signals. This thesis uses a combination of
machine-learning models and deterministic rules. The purpose is not to
replace managerial judgement, but to reduce the volume of raw data into a
smaller set of defensible indicators.

\subsection{Price forecasting with Prophet}

Price is a central competitive signal in online retail. It changes in
response to promotions, repricing strategies, inventory pressure, and
category-level competition. Classical autoregressive methods such as
ARIMA provide an important foundation for time-series forecasting, but
they often require careful stationarity assumptions and longer histories
than are available in short pilot deployments \citep{box2015time}.

Prophet, introduced by \citet{taylor2018forecasting}, is a decomposable
additive forecasting model designed for practical forecasting at scale.
It combines trend, seasonality, and holiday effects while tolerating
missing observations and abrupt trend changes. Although the current
thesis dataset is too short to exploit weekly or yearly seasonality
fully, Prophet remains appropriate as a transparent short-horizon
baseline for producing 7-day price forecasts and prediction intervals.
More complex forecasting models such as NeuralProphet and Temporal
Fusion Transformers provide useful future directions
\citep{triebe2021neuralprophet, lim2021temporal}, but they are not the
primary modelling choice in the implemented system because the thesis
prioritises interpretability and deployability under limited data.

\subsection{Review sentiment and aspect-level customer signals}

Customer reviews provide qualitative evidence that cannot be captured by
price or ranking alone. Sentiment analysis converts unstructured review
text into polarity labels or scores, allowing the system to track whether
customer opinion is improving, stable, or deteriorating. Earlier
sentiment methods relied on lexicons and rule-based models such as VADER
\citep{hutto2014vader}. Transformer-based language models later
improved contextual understanding through self-attention
\citep{vaswani2017attention}, BERT \citep{devlin2019bert}, and RoBERTa
\citep{liu2019roberta}.

This thesis uses the RoBERTa-based model
{\texttt{\seqsplit{cardiffnlp/twitter-roberta-base-sentiment-latest}}} through the
Hugging Face Transformers library \citep{wolf2020transformers,
loureiro2022timelms, cardiffnlp2026model}. Although the model was not
trained specifically on Amazon reviews, it provides a practical
pretrained baseline for classifying review sentiment in a thesis-scale
system. The review pipeline therefore treats model output as an
operational signal rather than as a perfect representation of customer
opinion.

Aspect-based sentiment analysis (ABSA) extends this idea by linking
sentiment to specific product features, such as battery life, sound
quality, build quality, or delivery experience \citep{hu2004mining,
pontiki2014semeval, zhang2022survey}. In the implemented project,
however, aspect information is not produced by a custom ABSA model.
Instead, the system reads Amazon's product-level aspect summary when it
is available. This distinction is important: the thesis uses aspect data
as a complementary marketplace signal, not as evidence that a new
aspect-extraction model was developed.

\subsection{Price-tier clustering with K-Means}

Competitive price movements are more meaningful when interpreted within
the correct market segment. A \$10 price drop may be dramatic in an
entry-level category but marginal in a premium segment. K-Means
clustering, originally formalised by \citet{macqueen1967kmeans}, is a
simple and interpretable unsupervised method for partitioning numerical
observations into clusters. Market segmentation literature similarly
emphasises the value of grouping products or customers into comparable
segments before drawing strategic conclusions \citep{wedel2000marketSegmentation}.

In the implemented system, K-Means is applied to product prices within
each category's daily Top-50 ranking. The resulting clusters are mapped
to entry, mid-range, and premium tiers according to their cluster
centres. This method is intentionally lightweight: it is not a full
demand-estimation model, but it provides a useful context for comparing
prices and identifying category-level price structure. The implementation
uses scikit-learn's K-Means and scaling utilities \citep{pedregosa2011sklearn}.

\subsection{Visual monitoring with perceptual hashing}

Product images play an important role in e-commerce conversion and brand
positioning. Competitors may update main images, add promotional badges,
change packaging visuals, or redesign listing creatives. Exact
byte-level comparison is not suitable for this task because trivial
compression changes can alter image files without changing visual
meaning. Perceptual hashing (pHash) addresses this problem by producing
a compact representation of visual structure, commonly using frequency
features such as the Discrete Cosine Transform (DCT) \citep{zauner2010phash}.

The project uses pHash to compare a product's current image with its
previous image hash, implemented with Pillow and the Python ImageHash
package \citep{pillow2026docs, buchner2024imagehash}. A Hamming-distance
threshold is then used to decide whether the difference is large enough
to indicate a meaningful visual change. This approach is computationally
inexpensive and appropriate for scheduled monitoring. Its limitation is
that it detects that an image has changed, but it does not explain
semantically what changed; semantic image interpretation is therefore
reserved for future work.

\section{Operational market metrics and rule-based intelligence}

In applied competitive intelligence, not every useful signal requires a
complex predictive model. Many operational decisions depend on composite
scores, threshold rules, and ranked-list changes. This thesis therefore
combines machine-learning outputs with deterministic market heuristics.

\subsection{Rank movement, entrant/exit detection, and sponsored visibility}

Best Sellers Rank (BSR) and Top-50 category rankings are central signals
for Amazon marketplace visibility \citep{amazon2024bsr}. A product
entering or exiting a category ranking can indicate competitive churn,
new product momentum, or temporary promotional effects. The implemented
entrant/exit module uses a set-difference comparison between today's and
yesterday's ranked ASINs. This method is simple but transparent, which is
valuable in a thesis context where each alert must be explainable.

Sponsored placements are another important source of competitive
visibility. In search advertising, share-of-voice measures how much
visibility a brand receives relative to competitors in a given query or
placement context \citep{ghose2009}. The implemented sponsored-share
module estimates brand-level sponsored visibility from
\texttt{category\_rankings.is\_sponsored}. This does not measure full
Amazon advertising spend; rather, it provides an observable proxy for
paid visibility within the collected category result pages.

\subsection{Composite scores: BMS and LQS}

Composite indicators are common in managerial decision support because
they condense several heterogeneous measurements into a single
interpretable score. The main risk of such scores is overclaiming:
weights and thresholds must be presented as transparent heuristics, not
as universally optimal parameters.

The Brand Momentum Score (BMS) in this thesis is a weighted operational
metric combining BSR velocity, review velocity, and sentiment. Its role
is to rank watchlist ASINs by recent competitive momentum, not to infer
causal drivers of sales. Similarly, the Listing Quality Score (LQS)
combines listing completeness and reputation-related features, including
title length, bullet count, image availability, A+ content availability,
star rating, review count, and sentiment. This design is consistent with
industry listing-quality tools \citep{amazon2024listing,
helium10listing, junglescout2024}, but the specific weights in the
implemented system should be interpreted as project-defined heuristics.

\subsection{Rule-based alerts}

Rule-based alerting remains useful when the goal is operational
interpretability. In this project, alerts are generated for events such
as price drops, stockouts, BSR deterioration, image changes, new
entrants, and exits from the Top 50. Each alert is tied to an explicit
condition or threshold, making the system easy to audit. This is
especially important because the downstream Telegram agents must provide
concise recommendations grounded in observable metrics rather than
unsupported language-model speculation.

\section{Evaluation under short-window deployment constraints}

Because this thesis is based on a short pilot deployment, the evaluation
strategy must be modest and transparent. The objective is to test whether
the implemented analytical components produce measurable and defensible
outputs, not to claim definitive long-term predictive superiority.

\subsection{Distant supervision for sentiment evaluation}

Manually labelled sentiment datasets are expensive to build. Distant
supervision offers a practical alternative by using naturally occurring
signals as noisy labels \citep{go2009twitter}. In review analysis,
star ratings can serve as a proxy for sentiment: one- and two-star
ratings are often treated as negative, three-star ratings as neutral, and
four- and five-star ratings as positive. This approach enables
large-scale evaluation, but it introduces label noise because review
text and star rating do not always express the same sentiment.

The implemented evaluation therefore reports not only accuracy, but also
per-class precision, recall, and macro-F\textsubscript{1}. Macro-F\textsubscript{1}
is particularly important because Amazon reviews are typically skewed
toward positive ratings. A classifier can appear accurate by predicting
the majority class, while still performing poorly on negative or neutral
reviews. Reporting macro-F\textsubscript{1} provides a more balanced view
of model behaviour under class imbalance.

\subsection{Walk-forward forecasting evaluation}

Time-series forecasting must be evaluated in a way that respects
temporal order. Random train-test splitting can leak future information
into the training set and therefore overstate performance. Walk-forward
or holdout-based backtesting avoids this problem by training on earlier
observations and evaluating on later observations.

The implemented forecast evaluation holds out the final three days for
eligible ASINs, fits Prophet on the earlier observations, and compares
predictions with actual held-out prices. It reports Mean Absolute Error
(MAE), Mean Absolute Percentage Error (MAPE), a naive last-observed-price
baseline, and 80\% prediction-interval coverage. The naive baseline is
especially important: if Prophet cannot outperform a simple
last-value-carry-forward forecast on short histories, the thesis must
disclose that limitation rather than overstate predictive value.

\section{Intelligence delivery: dashboards and tool-using agents}

Analytical outputs must be delivered in a form that supports decision
making. Dashboards and conversational agents serve different purposes,
and the implemented system uses both.

\subsection{Visual analytics and dashboards}

Interactive dashboards remain valuable for exploration and visual
analysis. Streamlit supports rapid development of data applications and
is appropriate for internal analytical dashboards \citep{streamlit2024docs}.
In this thesis, the Streamlit dashboard provides pages for market
overview, price intelligence, brand competition, listing quality,
sentiment and reviews, and alerts. A lightweight Vercel-deployable web
dashboard complements it by providing a compact public-facing overview.
The dashboards use Plotly for interactive charts \citep{plotly2024docs}.
Both dashboards are pull-based interfaces: they are useful when the user
chooses to inspect the system.

\subsection{Tool-using LLM agents and push-based delivery}

Recent work on tool-using language models shows that LLMs can become
more useful when grounded in external tools rather than relying only on
their internal parameters. Toolformer demonstrates language-model tool
use \citep{schick2023toolformer}, while ReAct shows how reasoning and
acting can be interleaved through tool calls \citep{yao2022react}.
Structured function-calling APIs further popularise this design pattern
\citep{openai2023functions}.

The present thesis applies this idea in a constrained and auditable way.
The OpenClaw layer does not allow agents to write to the database.
Instead, each agent calls read-only Python skills that query Supabase and
return JSON outputs. The three deployed agents are role-specific:
Competitor Spy focuses on rankings, BMS, sponsored share, entrant/exit
events, and image changes; Sentiment Detective focuses on sentiment,
reviews, aspects, and snapshots; Momentum Strategist synthesises alerts,
forecasts, BMS, LQS, and sentiment into action-oriented briefs. This
architecture reduces hallucination risk by grounding agent responses in
database-backed skill outputs, while also reducing dashboard friction by
pushing concise briefs through Telegram \citep{telegram2024bot}.

\section{Research gap and proposed contribution}

The literature and commercial landscape reveal four practical gaps.
First, robust scraping and data-extraction services exist, but they
typically provide raw datasets rather than thesis-specific analytical
signals. Second, machine-learning models for forecasting and sentiment
analysis are well established, but they are often evaluated separately
from the operational systems that deliver their outputs. Third,
commercial dashboards provide valuable market visibility, yet their
pull-based design still requires users to perform much of the monitoring
and synthesis manually. Finally, recent LLM-agent research demonstrates
tool use, but many applications remain generic and are not tightly
grounded in domain-specific, read-only analytical skills.

This thesis addresses these gaps by implementing a scheduled, low-cost,
multi-channel Amazon market-intelligence prototype. It integrates
Apify-based data acquisition, a Supabase PostgreSQL warehouse, Python
analytics scripts, interpretable ML components, composite operational
metrics, rule-based alerts, Streamlit and web dashboards, and an
OpenClaw Telegram delivery layer. The contribution is therefore not a
new standalone machine-learning algorithm; rather, it is the integration
and evaluation of existing data-engineering, machine-learning, and
tool-using-agent techniques in a working competitive-intelligence system
for a clearly bounded Amazon marketplace pilot.
